% SPDX-License-Identifier: CC-BY-4.0
%
% THERESA theory.  Build with:
%
%     latexmk -pdf docs/theory.tex -outdir=docs/_build
%
\documentclass[11pt,a4paper]{article}

\usepackage[T1]{fontenc}
\usepackage[utf8]{inputenc}
\usepackage{lmodern}
\usepackage{microtype}
\usepackage[a4paper,margin=28mm]{geometry}
\usepackage{amsmath,amssymb,amsthm}
\usepackage{booktabs}
\usepackage{siunitx}
\usepackage[hidelinks]{hyperref}
\usepackage{enumitem}

\theoremstyle{plain}
\newtheorem{proposition}{Proposition}

\sisetup{per-mode=symbol, range-phrase={\,--\,}, range-units=single}

% \year is a TeX primitive (the current year counter), so the unit needs its
% own name rather than redefining it.
\DeclareSIUnit{\annum}{a}

\newcommand{\Tss}{T_{\mathrm{ss}}}
\newcommand{\Ttop}{T_{\mathrm{top}}}
\newcommand{\Tlim}{T_{\mathrm{lim}}}
\newcommand{\dd}{\,\mathrm{d}}

\title{A surrogate model for the transient temperature field\\
       around an emplaced waste canister}
\author{THERESA 1.0.0 --- theory\thanks{\copyright\ 2026 frederik-f.
  Licensed under CC BY 4.0, \url{https://creativecommons.org/licenses/by/4.0/}.}}
\date{}

\begin{document}
\maketitle

\begin{abstract}
\noindent
A three-dimensional transient heat-conduction model of a Pollux-sized canister
in drift emplacement takes minutes per design. This note derives the
construction that reduces a parameter study to a single interpolation. Because
the governing equation is linear, the surface temperature, the geothermal
gradient, the emplacement depth, the canister loading and the decay-heat curve
--- its amplitudes, its decay constants and their number --- enter \emph{in
closed form}. The ten parameters that change the differential operator --- the
two spacings and the conductivity and volumetric heat capacity of four
materials --- are interpolated. The construction agrees with a conventional
simulation to \SI{0.07}{\kelvin} at the peak, the interpolation with unseen
designs to \SI{0.130}{\kelvin}.
\end{abstract}

\section{The model}

\subsection{Domain}

The repository is idealised as an infinite array of identical, simultaneously
emplaced canisters lying in parallel drifts. The array has four planes of
mirror symmetry through each canister, so it suffices to model one quarter of
the unit cell,
\[
  \Omega = \Bigl[0, \tfrac{p}{2}\Bigr] \times
           \Bigl[0, \tfrac{a}{2}\Bigr] \times
           \bigl[-H,\, 0\bigr],
\]
where $x$ runs along the drift axis, $y$ perpendicular to it, and $z$ is height
with $z = 0$ at ground level. The two lateral extents are half the canister
pitch $p$ and half the drift spacing $a$.

Neighbouring canisters are separated by a \emph{clear canister spacing} $g$,
measured between their end faces, from which the pitch follows:
\begin{equation}
  p = L + g,
  \label{eq:pitch}
\end{equation}
with $L$ the canister length. This matters for more than nomenclature: a pitch
would have to satisfy $p > L$, so treating $p$ and $L$ as independent design
variables would make the design space non-rectangular and generate impossible
designs in one corner. In $(g, L)$ the box is rectangular by construction.

$\Omega$ contains four subdomains with distinct, constant material properties:
the canister (a quarter cylinder of radius $r_{\mathrm{c}}$ and length $L/2$
within $\Omega$), the backfill, the liner, and the host rock. The drift runs the
full length of the cell; only the canister is finite.

\subsection{Governing problem}

With $\rho$, $c_p$ and $\lambda$ piecewise constant,
\begin{equation}
  \rho c_p \frac{\partial T}{\partial t}
    = \nabla\cdot\bigl(\lambda \nabla T\bigr) + q,
  \qquad \text{in } \Omega \times (0, t_{\mathrm{end}}],
  \label{eq:heat}
\end{equation}
subject to
\begin{equation}
  \begin{aligned}
    T &= \Ttop                  && \text{on } \Gamma_{\mathrm{top}} \ (z = 0),\\
    T &= \Ttop + \gamma H       && \text{on } \Gamma_{\mathrm{bot}} \ (z = -H),\\
    \lambda \nabla T \cdot n &= 0 && \text{on } \Gamma_{\mathrm{sym}},
  \end{aligned}
  \label{eq:bc}
\end{equation}
where $\gamma$ is the geothermal gradient and $\Gamma_{\mathrm{sym}}$ collects
the four lateral faces. A zero-flux condition is the natural boundary condition
of the weak form, so in a finite-element discretisation the symmetry planes are
imposed by \emph{not} constraining them.

The source is the canister's decay heat, distributed uniformly over the
canister volume:
\begin{equation}
  q(x,t) =
  \begin{cases}
    \dfrac{Q(t)}{4\,V_{\mathrm{c}}}, & x \in \Omega_{\mathrm{c}},\\[2ex]
    0, & \text{otherwise},
  \end{cases}
  \qquad
  Q(t) = \sum_{i=1}^{N} P_i\, e^{-\lambda_i t}.
  \label{eq:source}
\end{equation}
Here $Q$ is the power of the \emph{whole} canister and $V_{\mathrm{c}}$ the
volume of the quarter canister actually present in $\Omega$, so that
$\int_\Omega q \dd V = Q/4$: the quarter cell receives a quarter of the power,
as it must, since it represents a quarter of the unit cell. The factor is easy
to get wrong by four.

A second subtlety hides in $V_{\mathrm{c}}$. A tetrahedral mesh approximates
the cylinder from the inside, so the meshed canister is smaller than the
analytic one --- about \SI{0.45}{\percent} at the resolution used here. Since
the solver integrates $q$ over the \emph{meshed} elements, referring $q$ to the
analytic volume would inject that much too little heat, worth about a kelvin on
a rise of a few hundred. $V_{\mathrm{c}}$ is therefore the meshed volume.

\subsection{Initial condition}

The simulation starts from the undisturbed rock, i.e.\ from the source-free
steady state $\Tss$ solving
\begin{equation}
  \nabla\cdot\bigl(\lambda \nabla \Tss\bigr) = 0
  \label{eq:steady}
\end{equation}
with the same boundary conditions~\eqref{eq:bc}.

\section{Linearity}

Equation~\eqref{eq:heat} is linear in $T$: the coefficients $\rho$, $c_p$ and
$\lambda$ do not depend on $T$. Every reduction below follows from that single
fact, and every one of them fails the moment a temperature-dependent
conductivity is introduced.

\subsection{The steady state: surface temperature and gradient are free}

\begin{proposition}\label{prop:psi}
Let $\psi$ solve \eqref{eq:steady} with $\psi = 0$ on $\Gamma_{\mathrm{top}}$,
$\psi = H$ on $\Gamma_{\mathrm{bot}}$, and zero flux on
$\Gamma_{\mathrm{sym}}$. Then for every $\Ttop$ and every $\gamma$,
\begin{equation}
  \Tss(x) = \Ttop + \gamma\,\psi(x).
  \label{eq:tss}
\end{equation}
\end{proposition}

\begin{proof}
The right-hand side satisfies \eqref{eq:steady}, because $\Ttop$ is a constant
and $\nabla\cdot(\lambda\nabla(\Ttop + \gamma\psi)) = \gamma\,
\nabla\cdot(\lambda\nabla\psi) = 0$. It takes the value $\Ttop$ on
$\Gamma_{\mathrm{top}}$ and $\Ttop + \gamma H$ on $\Gamma_{\mathrm{bot}}$, and
its normal flux vanishes on $\Gamma_{\mathrm{sym}}$. The elliptic problem with
mixed Dirichlet--Neumann data has a unique solution, so the two coincide.
\end{proof}

$\psi$ depends on the geometry alone. One source-free solve per design
therefore covers \emph{every} surface temperature and \emph{every} geothermal
gradient exactly, at no further cost. In a homogeneous medium $\psi(x) = -z$;
the material contrast perturbs it only locally, by about
\SI{0.02}{\kelvin}-equivalent at the canister, and computing $\psi$ rather than
assuming a linear profile keeps that correction.

\subsection{The heating: the decay curve is free}

Write $u = T - \Tss$. Subtracting \eqref{eq:steady} from \eqref{eq:heat} and
using that $\Tss$ is time-independent,
\begin{equation}
  \rho c_p \frac{\partial u}{\partial t}
    = \nabla\cdot\bigl(\lambda \nabla u\bigr) + q,
  \qquad
  u(\cdot, 0) = 0,
  \label{eq:rise}
\end{equation}
with $u = 0$ on $\Gamma_{\mathrm{top}} \cup \Gamma_{\mathrm{bot}}$ and zero flux
on $\Gamma_{\mathrm{sym}}$. The boundary data of $u$ are homogeneous, so
\emph{$u$ does not depend on $\Ttop$ or $\gamma$ at all}. Together
with~\eqref{eq:tss},
\begin{equation}
  T(x,t) = \underbrace{\Ttop + \gamma\,\psi(x)}_{\text{one steady solve}}
         + \underbrace{u(x,t)}_{\text{one transient solve}} .
  \label{eq:split}
\end{equation}

Problem~\eqref{eq:rise} is linear and time-invariant in the source. Let
$S(x,t)$ denote its solution for the unit step $Q \equiv \SI{1}{\watt}$. Then
Duhamel's principle gives, for any power history $Q$,
\begin{equation}
  u(x,t) = \int_0^t \dot S(x, t - \tau)\, Q(\tau) \dd\tau
         = \int_0^t \dot S(x, v)\, Q(t - v) \dd v .
  \label{eq:duhamel}
\end{equation}

\subsubsection*{Evaluating \eqref{eq:duhamel} exactly}

$S$ is known at the solver's own time levels
$0 = t_0 < t_1 < \dots < t_{N_t}$. Between them the finite-element solution is
interpolated linearly in time, so taking $S$ piecewise linear is not an
approximation \emph{of the data} --- it is the data. Its derivative is then
piecewise constant,
\[
  \dot S(v) = s_j \equiv \frac{S_{j+1} - S_j}{t_{j+1} - t_j},
  \qquad v \in [t_j, t_{j+1}),
\]
and each interval integrates in closed form. For a single source term
$P e^{-\lambda t}$ and an output time $t_n$,
\[
  \int_{t_j}^{t_{j+1}} P\, e^{-\lambda (t_n - v)} \dd v
    = \frac{P}{\lambda}
      \Bigl( e^{-\lambda (t_n - t_{j+1})} - e^{-\lambda (t_n - t_j)} \Bigr),
\]
so that, summing over the terms and the intervals below $t_n$,
\begin{equation}
  \boxed{\;
  u(t_n) = \sum_{i=1}^{N} \frac{P_i}{\lambda_i}
           \sum_{j < n} s_j
           \Bigl( e^{-\lambda_i (t_n - t_{j+1})}
                - e^{-\lambda_i (t_n - t_j)} \Bigr). \;}
  \label{eq:closed}
\end{equation}
Every exponent in \eqref{eq:closed} is non-positive. The algebraically
equivalent form with $e^{+\lambda_i v}$ would overflow in double precision for
a short-lived nuclide over a \SI{1000}{\annum} window; this one cannot.

Two limiting cases are worth recording. For $\lambda_i \to 0$ the quotient tends
to the interval length $t_{j+1} - t_j$, and \eqref{eq:closed} collapses to
$u(t_n) = P\, S(t_n)$ --- a constant source reproduces the step response, which
is the identity that fixes the normalisation of $S$ and is checked to round-off
in the test suite. And if $S$ itself happens to be an exponential sum,
$S(t) = \sum_k a_k (1 - e^{-\mu_k t})$, term-by-term integration of
\eqref{eq:duhamel} gives the independent closed form
\begin{equation}
  u(t) = \sum_k \sum_i a_k \mu_k P_i \,
         \frac{e^{-\lambda_i t} - e^{-\mu_k t}}{\mu_k - \lambda_i},
  \label{eq:expkernel}
\end{equation}
with the removable singularity at $\mu_k = \lambda_i$ replaced by its limit
$a_k \mu_k P_i\, t\, e^{-\mu_k t}$. Equation~\eqref{eq:expkernel} is the
reference against which \eqref{eq:closed} is verified on a kernel whose exact
answer is known.

\subsubsection*{Consequence}

The amplitudes $P_i$, the decay constants $\lambda_i$ and the term count $N$
appear only in~\eqref{eq:closed}, never in a simulation. Six exponential terms
cost exactly what four or twenty cost, and introduce no error of their own. The
interim storage time $t_{\mathrm{s}}$ is likewise free: shifting an exponential
in time only rescales it,
\begin{equation}
  Q(t) = n \sum_i p_i e^{-\lambda_i (t + t_{\mathrm{s}})}
       = \sum_i \underbrace{n\, p_i\, e^{-\lambda_i t_{\mathrm{s}}}}_{P_i}
         e^{-\lambda_i t},
  \label{eq:storage}
\end{equation}
so it folds entirely into the amplitudes.

\subsection{The loading, forwards and backwards}

From \eqref{eq:rise}, scaling the source scales the response: $Q \mapsto \alpha
Q$ implies $u \mapsto \alpha u$. The loading --- as a package count $n$ in
\eqref{eq:storage}, or directly as a power --- is therefore a multiplication.

The inverse question is the one a design review asks, and it is equally exact.
Let $\Gamma_{\mathrm{c}}$ be the canister surface and $\Tlim$ a design limit.
Requiring $T \le \Tlim$ everywhere on $\Gamma_{\mathrm{c}}$ and at all times
gives, using \eqref{eq:split},
\[
  \Tss(x) + \alpha\, u(x,t) \le \Tlim
  \qquad \forall\, x \in \Gamma_{\mathrm{c}},\ t \in [0, t_{\mathrm{end}}],
\]
and since $u \ge 0$,
\begin{equation}
  \boxed{\;
  \alpha_{\max}
    = \min_{\substack{x \in \Gamma_{\mathrm{c}},\ t \\ u(x,t) > 0}}
      \frac{\Tlim - \Tss(x)}{u(x,t)} . \;}
  \label{eq:admissible}
\end{equation}
The largest admissible loading is $\alpha_{\max} Q(0)$. No search, no
iteration, no further simulation: one pass over the stored
$(x,t)$ grid. If $\Tss > \Tlim$ anywhere on the canister surface --- the
undisturbed rock is already too hot --- the constraint set is empty and no
loading works, which \eqref{eq:admissible} reports as $\alpha_{\max} \le 0$.

Note that the maximum over the surface is taken \emph{after} the convolution,
per time step, rather than maximising the step response first and convolving a
single curve.

The leading point on the surface does change with time --- for the reference
design it is the barrel at the mid-length plane for 292 of 296 time steps, and
one neighbouring point for the other four. What makes this harmless is that the
two leading points are never more than $1.3\times10^{-3}$ apart in relative
terms, so the two orders were measured to agree to better than
\SI{0.1}{\milli\kelvin}. Convolving first is therefore not a correction to a
present error but insurance: it costs nothing here, and it remains correct if
the range is ever widened to shapes where the leading points separate, where
maximising first would overestimate.

\subsection{Depth}

The emplacement depth $d$ acts twice. It sets the undisturbed temperature at
the canister, which is $\Ttop + \gamma d$ to within the local material
perturbation and is carried exactly by $\psi$ in~\eqref{eq:tss}. And it sets
the distance from the heated zone to the upper Dirichlet boundary, which can
only affect $u$.

The second effect is bounded by diffusion. With $\alpha_{\mathrm{th}} =
\lambda/(\rho c_p) \approx \SI{8.6e-7}{\metre\squared\per\second}$ for the host
rock, the thermal penetration depth after \SI{1000}{\annum} is
$\sqrt{\alpha_{\mathrm{th}} t} \approx \SI{180}{\metre}$, and at the time of the
peak only \SIrange{35}{55}{\metre}. At \SI{300}{\metre} of cover or more the
boundary is several penetration depths away.

This is an estimate, not a proof, so the model measures it: kernels computed at
$d \in \{300, 500, 700, 1000\}\,\si{\metre}$ with everything else held fixed
differ, after the analytic term is removed, by at most
\SI{0.062}{\kelvin} (Table~\ref{tab:errors}). Depth is therefore treated
analytically and is not interpolated. The vertical extent keeps a constant
clearance below the drift, so that the comparison isolates the upper boundary.

\section{Interpolation of the step response}

After Sections 2.1--2.4, the only quantities that still require simulation are
those that change the differential operator itself:
\[
  \pi = \bigl(g,\ a,\
    \{\lambda_m\}_{m=1}^{4},\ \{(\rho c)_m\}_{m=1}^{4}\bigr) \in \mathbb{R}^{10},
\]
the clear canister spacing and the drift spacing, which change the mesh, and
the thermal conductivity and volumetric heat capacity of the four materials,
which change the coefficients. The object to be learned is the map
$\pi \mapsto S(\cdot,\cdot;\pi)$.

Two remarks on this split. First, heat capacity appears only as the product
$\rho c$: the density and the specific heat never enter~\eqref{eq:heat}
separately, so treating them as two parameters would add a dimension carrying
no information. Second, the split is not between ``important'' and
``unimportant'' but between what the designer \emph{chooses} and what the
operator \emph{contains}. Everything exact in Sections 2.1--2.4 is a choice: a
surface temperature, a loading, a storage time. Everything interpolated is a
property. That the second group is also where the physical uncertainty lives is
what makes a variance-based sensitivity analysis over $\pi$ meaningful.

The canister length $L$ is deliberately \emph{not} in $\pi$. A canister type is
a discrete choice, and a length interpolated between two types corresponds to
no object; a different type is trained separately.

\subsection{Sampling the canister surface}

$S$ is needed on the canister surface, but the finite-element node sets of two
different designs are not comparable, so a nodal representation cannot be
interpolated across designs. Instead $S$ is sampled at $N_s = 156$ probe points
fixed in \emph{canister-local} coordinates: a grid in the circumferential angle
$\theta \in [0, \pi]$ and the normalised axial coordinate $\xi \in [0, 1]$ of
the half-length, plus rings on the end face. Because the grid is normalised, a
probe sits at the same relative position for every design, so the $N_s$
components of $S$ mean the same thing from one design to the next --- which is
what makes interpolation across designs well defined at all.

The sampling loss of that grid is measured against the exact nodal maximum over
the canister-wall submesh, and is recorded for every training run. It is
\SI{-0.01}{\percent}: the probe grid is marginally \emph{above} the nodal
maximum, because the probes sit \SI{2}{\centi\metre} inside the surface, where
the cast iron is a shade hotter.

\subsection{Reduced basis}

A kernel is an $N_t \times N_s$ array --- some $4.6 \times 10^4$ numbers --- so
fitting one regressor per entry is neither necessary nor sensible. Kernels of
different designs are very similar: they share a large near-field part, the
conduction drop through the backfill and liner immediately around the canister,
which depends on the spacings not at all and on $L$ only through $1/L$; and they
differ in a smaller, smooth array-interaction part.

A singular value decomposition of the mean-centred training kernels separates
exactly that. Writing $K(\pi) \in \mathbb{R}^{N_t \times N_s}$ flattened,
\begin{equation}
  K(\pi) \;\approx\; \bar K + \sum_{k=1}^{M} c_k(\pi)\, \Phi_k,
  \label{eq:pod}
\end{equation}
with $\Phi_k$ the leading right singular vectors. Modes are retained up to a
captured energy fraction of $1 - 10^{-8}$, at most $M = 8$; the remaining
modes carry run-to-run scatter only.

\subsection{Regression}

Each coefficient $c_k$ is a smooth scalar function of the ten variables, fitted
by a Gaussian process with a Matérn-$5/2$ covariance with one length scale per
input (ARD) and a white-noise term.

The inputs are transformed before fitting, conductivities and spacings
reciprocally and volumetric heat capacities directly:
\begin{equation}
  z = \bigl(1/g,\ 1/a,\ \{1/\lambda_m\}_{m=1}^{4},\ \{(\rho c)_m\}_{m=1}^{4}\bigr).
  \label{eq:transform}
\end{equation}
For the two spacings this reflects that neighbour interaction scales with
inverse distance. For the conductivities it reflects that a conduction
resistance is proportional to $1/\lambda$, so the response is close to linear
in $1/\lambda$ and far from it in $\lambda$. A volumetric heat capacity is not
a resistance --- it sets how fast the near field charges up --- and enters
roughly linearly, so it is sampled as it stands.

Training points are drawn from a Latin hypercube rather than a tensor grid. At
ten dimensions that is a necessity rather than a preference: six levels per
variable would be $6^{10} \approx 6 \times 10^{7}$ runs, whereas $n$ scattered
points give $n$ distinct levels in \emph{every} variable. Sampling is done in
the transformed coordinates~\eqref{eq:transform}.

\subsection{Uncertainty}

Each Gaussian process returns a posterior standard deviation $\sigma_k$ for its
coefficient. Every step from there to a temperature is linear ---
reconstruction~\eqref{eq:pod}, convolution~\eqref{eq:closed}, and the additive
split~\eqref{eq:split} --- so a one-sigma perturbation of $c_k$ maps to a
temperature curve $\delta T_k(t)$ by the same operations. The coefficients are
fitted independently, so the contributions combine in quadrature:
\begin{equation}
  \sigma_T(t) = \Bigl( \sum_{k=1}^{M} \delta T_k(t)^2 \Bigr)^{1/2},
  \qquad
  \delta T_k = \mathcal{D}\bigl[\sigma_k \Phi_k\bigr],
  \label{eq:sigma}
\end{equation}
where $\mathcal{D}$ denotes the convolution operator of \eqref{eq:closed}
evaluated at the probe point that is hottest at time $t$.

This is the uncertainty of the \emph{interpolation}, and of nothing else. It
says how unsure the regression is about a design it was not trained on. It does
not include the uncertainty of the physical model, which is dominated by the
material data and the modelling assumptions of Section~\ref{sec:limits} and is
very much larger.

\section{Discretisation and error budget}

\subsection{Time}

The schedule is blocked-logarithmic: a few steps at one size, then the size
grown by a fixed factor, until a ceiling is reached. Two things motivate this
rather than a smoothly growing schedule. The step response rises over six
decades of time and needs resolution at the fast end; and, since the problem is
linear, the solver may reuse its setup for as long as the step size is
unchanged, so keeping the number of \emph{distinct} sizes small is what makes
the run cheap. The shipped schedule has 295 steps but only 28 distinct sizes,
so \SI{9}{\percent} of steps pay for a solver setup.

\subsection{Space and solver}

Element size is graded from \SI{0.15}{\metre} at the canister to
\SI{30}{\metre} in the far field, reaching the coarse size only after
\SI{200}{\metre}. The last figure is not cosmetic: the field to be resolved
extends about \SI{180}{\metre}. The grading meshes every corner of the design
space.

Uniform refinement of this grading by $\sqrt{2}$ over four levels gives an
observed order of convergence of $2.22$ and a Richardson estimate
\SI{0.284}{\kelvin} above the shipped result, the approach being monotone from
below. The same exercise in time --- refining the whole schedule
self-similarly --- moves the peak by \SI{0.004}{\kelvin} with the sign
alternating, which is the solver's noise floor rather than a convergence rate.
The two are not comparable, and the shipped configuration is accordingly
generous in time and tight in space.

The linear systems are solved by conjugate gradients with diagonal
preconditioning. The system is symmetric positive definite --- the textbook
case --- and a direct sparse factorisation was measured to be at least
$3.7\times$ slower at this problem size.

\subsection{Budget}

\begin{table}[htbp]
  \centering
  \caption{Measured error contributions and their references. Accuracy target:
           \SI{0.5}{\kelvin} on the peak temperature.}
  \label{tab:errors}
  \begin{tabular}{llS[table-format=1.3]}
    \toprule
    Contribution & Measured against & {Error / \si{\kelvin}} \\
    \midrule
    Linear construction, Eq.~\eqref{eq:split}--\eqref{eq:closed}
      & conventional simulation & 0.070 \\
    Interpolation of the kernel
      & unseen simulated designs & 0.130 \\
    Treating depth analytically
      & kernels at four depths & 0.062 \\
    Time discretisation of \eqref{eq:closed}
      & analytic kernel, Eq.~\eqref{eq:expkernel} & 0.017 \\
    Probe-grid sampling
      & exact nodal maximum & 0.100 \\
    Spatial discretisation
      & Richardson extrapolation & 0.284 \\
    \bottomrule
  \end{tabular}
\end{table}


The first row is the decisive one. It compares the assembled answer --- one
steady solve, one step-response run, and the convolution with a six-term decay
curve --- against a conventional transient simulation of the same case with the
real decay curve as its source and the undisturbed field as its initial
condition. The two agree to \SI{0.070}{\kelvin} at the peak. Everything in
Section~2 rests on that agreement.

The last row is the largest, and it is of a different kind. Every other entry is
a difference between two things computed on the same mesh, so the mesh cancels
out of it; the spatial term is the distance from the mesh to the continuum, and
nothing in the other measurements can see it. It is reported here as a
\emph{documented offset} rather than applied as a correction. It is systematic
and one-signed --- all designs use the same mesh and all underestimate by about
the same amount --- so it largely cancels in comparisons between designs, which
is what the model is mostly used for. It does not cancel in an absolute value
compared against a temperature limit. Applying the extrapolated offset as a
constant would trade a known bias for an unknown one: it was measured at a
single design point, and there is no evidence that it is constant over a
ten-dimensional parameter space. An absolute peak needed to better than a few
tenths of a kelvin is best obtained from one directly simulated run on a refined
mesh at that design point.

\section{Assumptions and limits}\label{sec:limits}

\begin{itemize}[nosep,leftmargin=*]
  \item Pure heat conduction. No convection, no hydraulics, no mechanics, no
        radiative exchange across gaps, no resaturation of the backfill.
  \item Temperature-independent material properties. This is what makes the
        superposition exact, and it is the assumption that would have to go
        first if the model were extended: a conductivity $\lambda(T)$ invalidates
        Proposition~\ref{prop:psi}, Duhamel's principle~\eqref{eq:duhamel} and
        the loading relation~\eqref{eq:admissible} alike.
  \item Perfect thermal contact between all regions; no contact resistance.
  \item An infinite emplacement field with all canisters emplaced
        simultaneously --- the consequence of symmetry on all four lateral
        faces, and the conservative standard assumption for near-field design. A
        genuinely isolated canister would require far-field Dirichlet boundaries
        instead.
  \item The heat source is uniform over the canister volume.
  \item The host rock is homogeneous over the whole column: no geological
        layering, no depth-dependent conductivity or gradient.
  \item The interpolation is valid over the trained ranges of $g$, $a$ and $L$;
        outside them the prediction is an extrapolation and reports itself as
        one.
  \item THERESA is a private hobby project. Geometry, material parameters,
        decay-heat spectra and the design limit are simplified or assumed
        values, loosely based on published literature; they are not realistic
        and do not describe any existing or planned repository. The results
        are not suitable for design, licensing or safety assessments.
\end{itemize}

\end{document}
